

\cite{parhami1994multi}
defines impairments with coresponding views
component = defect
logic = fault
information = error
system = malfunction
service = degradation
result = failure

he mentions in each transition between the states, their are fields dedicated to - natural cause of the transition, impediment to the transition, techniques for faciliatting the transition, techniques for avoiding the transition, tools for modelling the transition

defect tolerance - redundacy


\cite{parhami2002defect}
failure is a matter of function, not whether something is fully intact or not. Failures can result from unsatifiable requirements
a system can be in one of seven states...: Ideal, Defective, Faulty, Erroneous, Malfunctioning, Degraded, Failed

Certain system states can expose a defect, resuling in a fault which is an incorrect signal values or decisions within the system. If the fault is exercises, it may contaiminte data flow in the system, causing errors. Erroneous information or states may or may not cause the affected subsystem to malfunction, depending on designed error tolerance. Degradation of service eventually leads to system failure.

A system may start up in any of the states. Moves between states through a series of deviations (events that take the systems to a lower state) and remedys (measures that enable the system to transition to a higher state)

Show example scenarios:
automobile break system

\cite{ferreira2023towards}
Define disturbance as "Something bad that may negatively affect the ability of a system (constituent or SoS) to deliver value or accomplish its mission."

define at SoS and Constituent level

Constituent disturbance - A disturbance that negatively affects a constituent system's ability to deliver value.
Classified by:
Origin - internal (component failure, software bug) or external (cyber attack, environmental event)

Intent - accidental (operator error) or malevolent (malicious input, attack)

SoS disturbance - A disturbance that negatively affects the SoS ability to accomplish its global mission
they don't have to result in failures, instead originate from 
interdependancies - 
"While interdependencies enable capabilities at the SoS level, they can also affect a SoS negatively. High levels of dependencies imply an increased risk of cascading effect."
"High levels of dependencies among constituents imply a higher risk of cascading failures in SoS."

Evolution of constituents - 
"Constituent systems can then exhibit unexpected behaviors over time."Not
"A deliberate change in a constituent system can produce severe consequences to SoS, generate cascading disturbances, and affect reliability at the SoS level."
"Failures in SoS may also occur due to the temporary inability of constituent systems to adequately provide functionalities."
"Such inability may not necessarily be related to failures in constituent systems rather to the constituent system's unavailability to meet the SoS demands at the appropriate time."


undesired emergent behavior -
"Undesired emergent behavior: Emergent behaviors that produce negative effects in SoS."
"As independent systems are connected to create a SoS, new behavioral properties can emerge from one or more systems behaving in unanticipated ways."


disturbance is better to use than fault as constituents may still be working in isolation, but behaviour in cooperation can negatively impact the SoS, therefore reliability degradation can still occur without a traditional fault or failure

\cite{ferreira2021reliability}
Insights from the literature review
The reliability of SoS requires up-to-date data. 
Models should represent SoS specific characteristics.
Optimized selection of constituent systems. 





\cite{axelsson2019refined}



\cite{baldwin2011formation}
belonging game matrix


\cite{salado2015abandonment}


\cite{salado2016exile}



% ----------------------------------------------------
Ferreira 2021 - Reliability in software-intensive systems: challenges, solutions, and future perspectives
Ferreira 2023 -  Towards an understanding of reliability of software-intensive systems-of-systems
Ferreira 2024 - A framework for the design of fault-tolerant systems-of-systems
Uday 2013 - Exploiting stand-in redundancy to improve resilience in a system-of-systems (SoS)
Imamura 2021 - System-of-systems reliability: An exploratory study in a brazilian public organization
Parhami 1994 - A multi-level view of dependable computing


WHERE THE HEALTH DIMENSION COMES FROM
\cite{ferreira2021reliability} in his SLR on relaibility in SoS he specifically mentions that factors affecting SoS relaibility are failures at the constituent level, communication problems, 

\cite{ferreira2023towards} defines diturbances across two acrosses about internal or external to the system and natural or accidental to the system, all these overall can be grouped into something that effects the ability of a system to delrivery what was intended and he outlines them by the causes:
"
Constituent System Disturbance Something bad that may negatively affect a constituent system’s ability to deliver value. It can be
classified along two axes; (1) internal or external to the system, or (2) natural/accidental or
malevolent.

Accidental disturbance A disturbance caused by unintentional actions (e.g., due to insufficient skills, stress, or fatigue). 

External disturbance A disturbance caused by system’s external factors (e.g., hacker attacks, natural disasters).

Internal disturbance A disturbance caused by system’s internal factors (e.g., components failures).

Malevolent disturbance A disturbance caused by malicious actions (e.g., erroneous inputs)
"


again in \cite{ferreira2024framework}, he mentions "In particular, a significant gap exists regarding the SoS dynamic nature. Constituent systems can produce disturbances at the SoS level that might result from failures, implementation changes, or, under certain circumstances, deviations from their roles within the SoS due to competing objectives. - all these lead to the idea of not being able to provide capabilities to the SoS as needed


- so this is where the idea of health originates from in the two dimensions where the idea if if unhealthy you are impedded from fulfilling the task you are meant to do. He just decribes the means this can occur in the paper

he also mentions \cite{avizienis2004basic}'s means to deals with faults and errors act as countermeasures to disturbances specifically in a later paper \cite{ferreira2024framework} he focuses on fault tolerance through mission decompisision in SoS. 

\cite{parhami1994multi-level} states align well with this as his progression in the reliability states is also centreded around not being able to deliver intended functionality



WHY MY APPROACH IS DIFFERENT
my approach brings the focus towards disturbances specifically 

In \cite{ferreira2023towards} he mkes a conceptial model of the relationships of concepts related to SoTS reliability derived from interviews with practitioners. His focus is on creating a mental model towards understanding rewlaibility, he mentios there no unique way to ensure relaibility due to the various issues and subfactors regarding the involved constituents


\cite{ferreira2023towards} mentions this in the discussion "This leads us to believe that understanding the impact of failures is as important as identifying their source...  In these cases, alternative actions should be taken" - I think this leads nicely to our solution in the case that were controlling eligibility, belonging, entry and exit in the SoTS rather than the constituents themselves. It provides an alternative to completely control over everything. 

\cite{ferreira2021reliability} he also specifically means a means to achieve reliability in SoS can be through model based analysis and specifcially \textbf{behaviour checking}, the system adheres to specific given properties  - I think this also leads to our solution of using the state charts to ensure consittuents follow certain behaviours depending on their state to prevent


SOTS TYPES 
\cite{ferreira2023towards} mentions this exact same idea but applied to countermeasures/fault tolerance to help improve SoS relaibility in the text too, which mainly comes around to having a centralized entity makes things easier. "A directed SoS, which is sometimes created to fulfill a specific mission and is centrally managed and coordinated to that end, is definable and potentially quantifiable in terms of reliability; so as in acknowledged SoS, although to a slightly lesser extent. On the other hand, the collaborative SoS, which counts on voluntary agreements of constituent systems, is likely to be much more challenging to adopt specific countermeasures.
